\documentclass[11pt]{article}
\usepackage{amsmath,amssymb,amsthm}
\usepackage[margin=1.1in]{geometry}
\usepackage{hyperref}

\newtheorem{proposition}{Proposition}
\newtheorem{remark}{Remark}
\newcommand{\sigmoid}{\sigma}
\newcommand{\xbar}{\bar{x}}

\title{Bernoulli Maximum-Likelihood Estimation, Part II:\\
The Practical Solution --- Turning Estimation into Optimization}
\author{Yuval Lavie}
\date{\today}

\begin{document}
\maketitle

\begin{abstract}
Part~I (\texttt{mle\_theoretical.tex}) solved the Bernoulli estimation
problem the statistician's way: maximize the likelihood, in closed form.
This part performs the machine-learning translation --- likelihood
$\to$ loss --- and solves the \emph{same} problem the way we must when no
closed form exists: gradient descent with a hand-derived gradient,
stochastic gradient descent from one sample at a time, and SGD with the
gradient computed by PyTorch's autograd. Two companion scripts share
Part~I's dataset: \texttt{mle\_practical\_pure.py} implements GD and SGD
by hand and \emph{animates} the estimator descending the loss landscape
alongside its loss progression; \texttt{mle\_practical\_pytorch.py}
drives autograd through the identical sample schedule and verifies the
trajectories coincide to machine precision --- a proof by computation
that autograd performs exactly our chain rule.
\end{abstract}

\section{From likelihood to loss}

Part~I ended with the log-likelihood
$\ell(p) = m\log p + (n-m)\log(1-p)$ and its maximizer
$\hat{p} = m/n = \xbar$. Machine learning prefers to \emph{minimize
losses}, and two cosmetic transformations get us there --- neither moves
the optimum. Dividing by the constant $n > 0$ rescales the function but
does not move its argmax (and makes the scale independent of dataset
size). Negating flips maximization into minimization. With
$\xbar = m/n$, the resulting \emph{average negative log-likelihood}
(NLL) is
\begin{equation}
  \mathcal{L}(p) \;=\; -\tfrac1n\,\ell(p)
  \;=\; -\xbar\log p - (1-\xbar)\log(1-p),
  \label{eq:nll}
\end{equation}
and the chain of equivalences
\begin{equation}
  \arg\max_p L(p)
  \;=\; \arg\max_p \ell(p)
  \;=\; \arg\min_p \mathcal{L}(p)
  \label{eq:equiv}
\end{equation}
says minimizing this loss \emph{is} maximum likelihood. Written as an
average of per-sample terms,
\begin{equation}
  \mathcal{L}(p) \;=\; \frac1n \sum_{i=1}^{n} \mathcal{L}_i(p),
  \qquad
  \mathcal{L}_i(p) \;=\; -x_i \log p - (1-x_i)\log(1-p),
  \label{eq:persample}
\end{equation}
which is the \emph{cross-entropy} loss --- and the sum-of-per-sample
structure is precisely what SGD will exploit.

\section{One coordinate change --- not a new objective}

Gradient methods want an unconstrained variable, but $p$ lives in
$(0,1)$. So we evaluate the \emph{same} NLL \eqref{eq:nll} at
\begin{equation}
  p = \sigmoid(t) = \frac{1}{1+e^{-t}}, \qquad t \in \mathbb{R},
  \qquad
  F(t) \;=\; \mathcal{L}\bigl(\sigmoid(t)\bigr).
  \label{eq:reparam}
\end{equation}
Nothing about the target function changes; we have only composed it with
a smooth bijection of the search space. (By the \emph{invariance} of
maximum likelihood --- Part~I's closing remark --- the optimum maps
back: $\hat p = \sigmoid(\hat t\,)$.)

\paragraph{The gradient, by the chain rule.}
Using $\sigmoid'(t) = \sigmoid(t)\bigl(1-\sigmoid(t)\bigr)$ and writing
$s = \sigmoid(t)$ for the moment,
\begin{equation*}
  F'(t)
  \;=\; \mathcal{L}'(s)\,\sigmoid'(t)
  \;=\; \left[-\frac{\xbar}{s} + \frac{1-\xbar}{1-s}\right] s(1-s)
  \;=\; -\xbar(1-s) + (1-\xbar)\,s
  \;=\; s - \xbar ,
\end{equation*}
hence
\begin{equation}
  \boxed{\;F'(t) \;=\; \sigmoid(t) - \xbar\;}
  \label{eq:grad}
\end{equation}
Every term cancelled except the residual: current estimate minus sample
mean. Moreover $F''(t) = \sigmoid'(t) > 0$, so $F$ is convex in $t$ ---
one basin.

\section{Route 1: gradient descent}

GD iterates the full gradient, every step:
\begin{equation}
  t \;\leftarrow\; t - \eta\,\bigl(\sigmoid(t) - \xbar\bigr).
  \label{eq:gd}
\end{equation}
This is deterministic: with a sane step size it converges monotonically
to the unique fixed point $\sigmoid(t^{\ast}) = \xbar$ --- Part~I's
closed form, now reached by iteration. For this toy problem the ``full
gradient'' is free ($\xbar$ is precomputed); in real models it costs a
pass over the entire dataset per step, which is what motivates the next
route.

\section{Route 2: stochastic gradient descent}

SGD replaces the full gradient \eqref{eq:grad} with a \emph{one-sample
estimate}: draw an index $i$ and use the gradient of the per-sample term
\eqref{eq:persample} through the same coordinate change,
\begin{equation}
  g_i(t) \;=\; \sigmoid(t) - x_i .
  \label{eq:sgdgrad}
\end{equation}
It is unbiased --- averaging over a uniformly drawn $i$,
$\mathbb{E}_i\!\left[g_i(t)\right] = \sigmoid(t) - \xbar = F'(t)$ --- so
each cheap step points in the right direction \emph{on average}. The
script sweeps reshuffled passes over the data, updating
$t \leftarrow t - \eta\, g_i(t)$ one sample at a time.

\paragraph{The price of stochasticity, and Polyak--Ruppert averaging.}
With a constant step size the SGD iterate never converges: it reaches a
stationary \emph{distribution}, a noise ball around the optimum whose
radius grows with $\eta$ and the per-sample variance. The estimator is
therefore taken as the average of the final pass's iterates,
$\hat{p}_{\mathrm{SGD}} = \tfrac1K \sum_k \sigmoid(t_k)$, which cancels
the hovering and lands within a few $10^{-3}$ of $\xbar$.

\begin{remark}[Mini-batches, later]
Averaging $g_i$ over a small batch $B$ gives $\sigmoid(t)-\xbar_B$: the
same expectation with variance reduced by $|B|$. Everything in this note
goes through unchanged; batching is an engineering trade between noise
and cost per step, not new mathematics.
\end{remark}

\section{Route 3: SGD with autograd}

PyTorch's \texttt{binary\_cross\_entropy\_with\_logits}$(t, x_i)$
evaluates exactly the one-sample NLL at $p = \sigmoid(t)$, i.e.\ the
composition $F_i(t) = \mathcal{L}_i(\sigmoid(t))$ of
\eqref{eq:persample} with \eqref{eq:reparam}. Calling
\texttt{loss.backward()} applies reverse-mode automatic differentiation
--- a mechanized chain rule --- so it must produce \texttt{t.grad}
$= \sigmoid(t) - x_i$, identically \eqref{eq:sgdgrad}.
\texttt{torch.optim.SGD} then performs the same update
$t \leftarrow t - \eta\,\texttt{t.grad}$. Autograd changes who does the
calculus, not the mathematics.

\paragraph{Proof by computation.}
The script drives the hand-written SGD and the autograd SGD through the
\emph{identical} sample schedule, the same $\eta$, and the same
initialization $t_0 = 0$ (in \texttt{float64} on both sides). Under
those conditions the two trajectories can differ only if the gradients
differ; the script asserts
$\max_k |\sigmoid(t_k^{\text{hand}}) - \sigmoid(t_k^{\text{autograd}})|
< 10^{-10}$, and measures the gap to be exactly $0.0$ over all
$15{,}001$ iterates --- an end-to-end check that \eqref{eq:sgdgrad} is
precisely what autograd computes.

\section{Experiment}

With $p = 0.3$ and $n = 5000$ (seed 7, Part~I's dataset):
GD ($\eta = 1$, $150$ steps) lands $5.6\times10^{-17}$ from the closed
form; SGD ($\eta = 0.1$, $3$ epochs), Polyak-averaged over the last
pass, lands within $1.5\times10^{-4}$; and the two SGD trajectories
coincide exactly.

\texttt{mle\_practical\_pure.py} plays the animation, three synchronized
views of the same descent: the loss landscape \eqref{eq:nll} with GD's
bead gliding down while SGD's bead jitters into its noise ball; the
estimate-versus-time curves approaching the closed-form line; and the
\emph{loss progression} --- the suboptimality gap
$\mathcal{L}(p_k) - \mathcal{L}(\hat p)$ on a log scale, where GD's
geometric convergence is a straight plunge and SGD's noise floor is a
plateau. \texttt{mle\_practical\_pytorch.py} animates the proof by
computation: the hand and autograd trajectories drawing simultaneously
and never separating, beside their per-step gap pinned to zero on a log
scale.

Three routes, one destination: the loss surface --- not the optimizer
--- decides where estimation lands; the optimizer only decides how you
travel and what noise you accept along the way.

\end{document}
